\cvsection{Expériences professionnelles}

\begin{cventries}
  \cventry
    {Stagiaire ingénieur de recherche en IA}
    {INRIA, équipe Mnémosyne}
    {Bordeaux, France}
    {Mars 2026 -- Septembre 2026}
    {
      \begin{cvitems}
        \item {Conception et développement en autonomie, avec le suivi de mes encadrants, d’un agent IA multi-étapes répondant au besoin d’aider les chercheurs à trouver des articles scientifiques.}
        \item {Construction de l’agent sans framework d’orchestration afin de comprendre et maîtriser l’exécution des outils, l’état, la mémoire, le contexte et les enchaînements de tâches.}
        \item {Intégration de l’API OpenAI et conception des outils utilisés par l’agent pour rechercher des informations et exécuter des tâches successives.}
        \item {Développement de services Python/FastAPI, d’API REST documentées avec OpenAPI et de la persistance MongoDB.}
        \item {Mise en œuvre d’un pipeline RAG comprenant le chunking, les embeddings, la recherche sémantique et une base vectorielle ChromaDB.}
        \item {Conteneurisation de composants avec Docker et mise en place d’un pipeline CI/CD sur GitHub.}
        \item {Gestion d’une application web conteneurisée permettant d’utiliser des LLM installés localement sur un Mac Studio, incluant l’installation et la configuration des modèles.}
        \item {Évaluation d’OpenClaw dans un environnement Docker isolé avec une attention portée à la fiabilité et à la sécurité des accès autonomes aux outils et aux fichiers.}
        \item {Comparaison de Qwen, Gemma et GLM et de plusieurs configurations de quantification sur une infrastructure HPC avec Slurm selon la latence, la TTFT, la qualité et la taille du contexte.}
        \item {Définition du protocole expérimental, analyse des résultats et contribution à un article scientifique sur l’évaluation des modèles.}
      \end{cvitems}
    }

  \cventry
    {Stagiaire de recherche en apprentissage par renforcement}
    {Université Karunya}
    {Coimbatore, Tamil Nadu, Inde}
    {Juin 2025 -- Sept. 2025}
    {
      \begin{cvitems}
        \item {Conception et comparaison d’une planification de trajectoire déterministe et d’une approche fondée sur l’apprentissage par renforcement pour la navigation autonome d’un robot hospitalier.}
        \item {Déploiement et optimisation du logiciel et du modèle d’apprentissage par renforcement sur un Raspberry Pi embarqué.}
      \end{cvitems}
    }
\end{cventries}
